\section{Introduction}
\label{sec:intro}

Autoregressive language models generate text one token at a time;
each token requires a forward pass over the growing KV cache. This is
the decoding equivalent of a toll road: every token pays a fare, and
the fare never decreases. Block-diffusion models generate an entire
block of tokens in parallel through iterative denoising, then step to
the next block. The cost is a turnstile: pay once, pass through a
whole block. We call this second class \emph{block-diffusion} because
the block structure --- not the diffusion process --- is what
determines its cost shape.

We compare the two on a pinned serving stack: vLLM on an NVIDIA A100
80GB, BF16. The task is structured log extraction from the LOGHUB BGL
corpus --- system logs from the BlueGene/L supercomputer at Lawrence
Livermore National Laboratory. Each raw log line describes a hardware
or software event; the model must extract a structured JSON record
with twelve fields spanning structural, semantic, anomaly, and action
categories.\footnote{Structural fields require span-copying from the
input; semantic fields require inference; anomaly fields require
binary classification. Quality is macro-averaged token-level F1 across
9 of the 12 fields; the other 3 have degenerate gold labels on BGL.
All experiments use $n = 300$ test events per cell unless stated
otherwise.}

Both arms are Mixture-of-Experts models:
\begin{itemize}[nosep]
\item \DG\ (DiffusionGemma-26B-A4B, 4.0\,B active of 26\,B total,
  52\,GB) --- the block-diffusion arm. Top-8 of 128 experts.
\item \ARthree\ (Qwen3-30B-A3B, 3.0\,B active of 30\,B total,
  60\,GB) --- the autoregressive arm. Matched to DG on architecture
  class and active parameters.
\end{itemize}
The cost shapes we derive are properties of the \emph{decoding
paradigm}, not of model scale. A larger AR model would have a steeper
slope but the same linear shape; a larger diffusion model would have a
higher floor but the same staircase.

Our contributions are as follows:
\begin{itemize}[leftmargin=1.5em]
\item \textbf{How do the two architectures differ in latency cost?}
We derive two cost laws from measurement: AR latency is a line in output
length with no floor; block-diffusion latency is a staircase, flat inside
a fixed-size block and stepping at each boundary. The shapes cross at
$L^{*} = \crossoverL$ output tokens (\S\ref{sec:laws}).

\item \textbf{Where is the diffusion model faster?}
Above $L^{*}$, DG is faster --- fifteen times at natural output, eight
times at a thousand tokens. In cost terms, that is the difference
between \$0.0076 and \$0.0005 per event on an A100, a gap that
compounds to thousands of dollars at million-event scale
(\S\ref{sec:wins}).

\item \textbf{What limits the diffusion model's advantage in practice?}
The AR model batches far better and has a tighter tail. But in the
batch-1 regime --- serial chains, autonomous driving, trading ---
batching is unavailable and the staircase shape dominates
(\S\ref{sec:limits}).

\item \textbf{When is block diffusion the right choice?}
When output is long, batch-1 is required, and schema completeness
matters. We state five conditions and ground them in concrete
workloads (\S\ref{sec:survives}).
\end{itemize}

\begin{table}[h]
\centering
\caption{The paper in one table. The diffusion model's latency
advantage is real; deployment factors limit it in practice.}
\label{tab:summary}
\begin{tabularx}{\linewidth}{@{}L r r@{}}
\toprule
& \thead{\ARthree\,(MoE)} & \thead{\DG-26B-A4B\,(MoE)} \\
\midrule
\tabfont Active parameters                            & 3.0\,B & 4.0\,B \\
\tabfont Total parameters                             & 30\,B & 26\,B \\
\tabfont Weights (bf16)                                & 60\,GB & 52\,GB \\
\tabfont Cost law slope (ms/token, serial) & 47.7 & --- \\
\tabfont Batch-1 $\to$ peak throughput multiplier     & 77$\times$ & $\leq$1.7$\times$ \\
\tabfont Tail latency (max/median)                    & 1.03$\times$ & 3.21$\times$ \\
\bottomrule
\end{tabularx}
\end{table}
